\section{Matter: the Kohn-Sham analysis}
\label{sec:ks}

This section determines the states of uniform matter made of the fermions of the model when the
self-interaction acts between them. The Vorton-Scalar is set aside here ($\yVS=0$, so that the mass
$m$ of the fermions in vacuum is the bare Dirac mass); Sec.~\ref{sec:cjtcontact} restores it.

\subsection{The method}
\label{sec:ksmethod}

In the method of Kohn and Sham \cite{KohnSham1965} a system of interacting fermions is replaced by
independent fermions that move in fields which their own densities determine. The energy is written
as a function of the densities; the fields are its derivatives; and the densities are computed from
the occupied states in those fields, until both agree. For uniform matter of the present model the
densities are two numbers (Sec.~\ref{sec:hf}): the number density $n=\langle\psi^\dagger\psi\rangle$
and the scalar density $\nS=\langle\bar\psi\psi\rangle$. The interaction energy is
$\tfrac12\GV n^{2}+\tfrac12\GS\nS^{2}$ with a scalar coupling $\GS$ and a vector coupling $\GV$, and
each particle obeys the Dirac equation with the Hamiltonian matrix
\begin{equation}
\begin{aligned}
\hh_{\mathrm{KS}}&=\sum_j\alpha^{j}k_j+\beta\,\mst+V_0,\\
\mst&=m+\GS\,\nS,\qquad V_0=\GV\,n .
\end{aligned}
\label{eq:hKS}
\end{equation}
The scalar field changes the mass of the particles in matter from $m$ to $\mst$; the vector field
shifts all energies by $V_0$. A negative $\GS$ is an attraction (it lowers the mass and the energy)
and a positive $\GV$ a repulsion; this is the structure of the mean-field theory of dense matter
with a scalar and a vector field \cite{Walecka1974}. With the couplings of Eq.~(\ref{eq:GHF}) the exchange terms are
contained exactly, because for a contact interaction they are a function of the same two densities.

A particle of the momentum $\vec k$ has the energy $E^{*}(k)+V_0$ with
$E^{*}(k)=(k^{2}+\mstq)^{1/2}$. The occupied states are those with $|\vec k|<\kF$, eight for each
momentum ($\nu=8$), and averages over them are written
$\langle f\rangle=\nu\int d^{D}k\,f/(2\pi)^{D}$, where $D$ is the number of space dimensions in which
the particles have momentum. Then
\begin{equation}
n=\langle1\rangle=\frac{\nu\,V_D\,\kF^{D}}{(2\pi)^{D}},\qquad\nS=\Big\langle\frac{\mst}{E^{*}}\Big\rangle,
\label{eq:densities}
\end{equation}
with the volume $V_D=\pi^{D/2}/\Gamma_{\mathrm E}(D/2+1)$ of the unit ball ($\Gamma_{\mathrm E}$ is
Euler's gamma function). The second of Eqs.~(\ref{eq:hKS}) with the second of
Eqs.~(\ref{eq:densities}) is the \emph{gap equation} for $\mst$. The energy density is
\begin{equation}
\varepsilon=\langle E^{*}\rangle+\tfrac12\GV\,n^{2}-\frac{(\mst-m)^{2}}{2\GS},
\label{eq:epsKS}
\end{equation}
which is the kinetic and mass energy of the particles plus the interaction energy [the first term
counts the mass as $\mst$ instead of $m$, which the last term corrects]; it is stationary with
respect to $\mst$ exactly when the gap equation holds. The pressure along the coordinate $xj$ is
\begin{equation}
p_j=\Big\langle\frac{k_j^{2}}{E^{*}}\Big\rangle+\pint,\qquad\pint=\tfrac12\GV\,n^{2}+\frac{(\mst-m)^{2}}{2\GS}.
\label{eq:pKS}
\end{equation}
The first term is the momentum that the particles carry across a surface; the \emph{interaction
pressure} $\pint$ is the same in all directions. \emph{Along a coordinate on which the fields do not
depend the particles have no momentum, and the pressure is $\pint$ alone}; this remark is used for
the deflating coordinates below. The chemical potential is $\mu=E^{*}(\kF)+V_0$, and for an
isotropic state $p=\mu n-\varepsilon$. Matter is \emph{self-bound} if its energy per particle,
$\varepsilon/n$, is below the mass $m$ at some density; the density at which $\varepsilon/n$ is
smallest is the \emph{saturation} density, where the pressure vanishes and
$\mu=\varepsilon/n$ \cite{HugenholtzVanHove1958}.

Two limitations are stated here once. The filled states of negative energy are not counted (the
\emph{no-sea approximation}); Sec.~\ref{sec:cjtcontact} counts them. And the state is the one of
lowest energy among the solutions of the gap equation: without the filled negative-energy states the
energy as a function of an unrestricted trial mass has no lowest value, because a large negative
trial mass places particles in states of negative rest energy.

%%<1671>
\subsection{Flat spacetime: the reference calculation}
\label{sec:ksflat}

On the flat $7+1$ spacetime the particles have momentum along all seven space coordinates, $D=7$,
and a coupling has the dimension of a length to the sixth power. The reference example has
$\GS=-10^{5}$ and $\GV=+10^{4}$ in units $m=1$. Table~\ref{tab:ksflat} gives its saturation point,
computed by the program of this paper, together with the values of a first calculation with a
notebook in the Wolfram Language; an independent script reproduces the same numbers a third time.
Figure~\ref{fig:ksflat} shows the equation of state.

\begin{table*}
\caption{\label{tab:ksflat}Uniform matter on the flat $7+1$ spacetime at saturation, for
$\GS=-10^{5}$, $\GV=+10^{4}$ (units $m=1$, seven space dimensions): the program of this paper and
the first calculation with a notebook.}
\begin{ruledtabular}
\begin{tabular}{lll}
quantity & this program & notebook \\
\hline
{{ROWS:tbl_ks_flat}}
\end{tabular}
\end{ruledtabular}
\end{table*}

\begin{figure*}
\includegraphics[width=\textwidth]{fig_ks_flat}
\caption{\label{fig:ksflat}Uniform matter on the flat $7+1$ spacetime (seven space dimensions; units
$m=1$) for two pairs of mean-field couplings. (a) Energy per particle minus $m$; the circle marks the
saturation point of Table~\ref{tab:ksflat}. (b) Mass of the particles in matter. (c) Pressure
divided by the energy density.}
\end{figure*}

The matter is self-bound: at $\kF={{GENM:ks_kf_sat}}\,m$ the energy per particle is
${{GENM:ks_ebind}}\,m$ below the mass, and the mass of the particles in it is
$\mst={{GENM:ks_mstar}}\,m$. The identity $p=\mu n-\varepsilon$ holds along the whole equation of state
to $10^{-15}$ of $\varepsilon$, and at saturation $\mu=\varepsilon/n$.

One remark concerns the relation between these two numbers and the couplings of the Lagrangian. If
the Lagrangian has $g_0=-10^{5}$, $g_1=+10^{4}$ and no other coupling, Eq.~(\ref{eq:GHF}) gives
$\GS^{\mathrm{HF}}=-98750$ and $\GV^{\mathrm{HF}}=+20000$: the exchange terms double the vector
repulsion. For that pair the energy per particle is above $m$ at all densities [dashed curves of
Fig.~\ref{fig:ksflat}; its local minimum lies ${{GENM:ks_hf_localmin}}\,m$ above $m$, at
$\kF={{GENM:ks_hf_localmin_kf}}\,m$]: \emph{uniform matter is not self-bound.} The reference example is therefore to be read as $\GS=-10^{5}$, $\GV=+10^{4}$ for the
mean-field couplings themselves, which corresponds to $g_0=-105341.6$, $g_1=+2484.5$ if no other
coupling of the Lagrangian differs from zero.

\subsection{Matter with momentum along the deflating coordinates}
\label{sec:ks71}

Now let $x5,x6,x7$ deflate, and consider matter whose particles have momentum along all seven space
coordinates. The linear Dirac equation conserves the wavenumbers $k_t$ of a particle
(Sec.~\ref{sec:dirac}); its proper momentum along $xt$ is $k_t/b_t$ and grows as the coordinate
shrinks. If the particles do not collide, the occupied states remain the ball $|\vec k|<\kF$ of
wavenumbers that they filled at a reference time, and the occupied proper momenta form an
\emph{ellipsoid} that is stretched along $x5,x6,x7$ by the factors
$b_t^{\mathrm{ref}}/b_t=e^{\rt(x4-x4_{\mathrm{ref}})}$. The number density grows as the volume
shrinks,
\begin{equation}
n=n_{\mathrm{ref}}\,e^{\td},
\end{equation}
and the equations of Sec.~\ref{sec:ksmethod} hold with the averages taken over the ellipsoid; the
pressure along a deflating coordinate is $p_t=\langle(k_t/b_t)^{2}/E^{*}\rangle+\pint$. The energy
density obeys
\begin{equation}
\frac{d\varepsilon}{dx4}=\sum_{t}\rt\,(\varepsilon+p_t),
\label{eq:conservation}
\end{equation}
the statement that the covariant divergence of the energy-momentum tensor vanishes; the program
confirms it to $6\times10^{-9}$ and the independent script to $2\times10^{-8}$.

Figure~\ref{fig:ksdefl} shows the history of the saturated matter of Table~\ref{tab:ksflat}, taken as
the state at the reference time, for three equal rates and for rates in the ratio $1:2:3$.

\begin{figure*}
\includegraphics[width=\textwidth]{fig_ks_deflating}
\caption{\label{fig:ksdefl}Matter of the $7+1$ model with momentum along all seven space coordinates
while $x5,x6,x7$ deflate, without collisions. At the clock reading $\td=0$ it is the saturated matter
of Table~\ref{tab:ksflat}. (a) Mass of the particles in matter. (b) Energy per particle minus $m$:
for the collisionless state, and for the state of lowest energy at the same density (``relaxed'').
(c) Pressures along the flat coordinates and along the three deflating ones, divided by the energy
density, for rates in the ratio $1:2:3$.}
\end{figure*}

\emph{The mass of the particles falls.} It reaches half its value at the reference time at
$\td={{GENM:ks71_half_eq}}$ (equal rates) or ${{GENM:ks71_half_123}}$ (rates $1:2:3$).

\emph{Matter is self-bound only for a limited time.} The energy per particle is below $m$ only for
${{GENM:ks71_lo_eq}}<\td<{{GENM:ks71_hi_eq}}$ (equal rates;
${{GENM:ks71_lo_123}}<\td<{{GENM:ks71_hi_123}}$ for $1:2:3$): about one e-folding of the deflating
volume. Before that the matter is too dilute to bind, afterwards the vector repulsion
$\tfrac12\GV n^{2}$ dominates; at $\td=4$ it is {{GEN:ks71_vec_eq}} per cent of the energy density,
which then grows as $e^{2\td}$.

\emph{The pressure becomes anisotropic}, larger along the coordinates that deflate faster
[Fig.~\ref{fig:ksdefl}(c)].

\emph{Collisions would release heat.} The collisionless state has more energy than the state of
lowest energy at the same density, a ball of occupied states; the difference reaches
${{GENM:ks71_heat}}\,m$ per particle at $\td=4$. Whether collisions are frequent enough to release it
is not examined here.
%%</1671>

\subsection{Fields that do not depend on the deflating coordinates}
\label{sec:ksreduced}

%%<1644>
On the $4+4$ spacetime matter exists only in the quantized sector: a mode that depends on a
deflating time is extinguished (Sec.~\ref{sec:modes}). The particles then have momentum along
$x0,\dots,x3$, so that $D=4$, or, if the fields do not depend on the hidden coordinate $x0$ either,
along ordinary space, $D=3$. The self-interaction is taken in the quantized-sector form
(Sec.~\ref{sec:forms}).
%%</1644>
%%<1671>
Consider now the sector of fields that do not depend on $x5,x6,x7$. Its particles have momentum
along $x0,\dots,x3$, so that $D=4$, or, if the fields do not depend on the hidden coordinate $x0$
either, along ordinary space, $D=3$.
%%</1671>
By Eq.~(\ref{eq:couplings}) the equations of Sec.~\ref{sec:ksmethod} hold for the weighted field at
each moment with the couplings of that moment,
\begin{equation}
\GS(x4)=\GS^{\mathrm{ref}}\,e^{\td},\qquad\GV(x4)=\GV^{\mathrm{ref}}\,e^{\td},
\end{equation}
while the mass $m$ and the density of particles per unit of the volume of $x0,\dots,x3$ do not
change. The ratio $\GV/|\GS|$ is fixed. The deflation therefore carries the matter through a
one-parameter sequence of interactions of growing strength, and \emph{a gas becomes self-bound
matter at a definite moment}: when the scalar attraction reaches the threshold of
Table~\ref{tab:ksthr}.

\begin{table}
\caption{\label{tab:ksthr}The smallest scalar attraction at which uniform matter is self-bound, for
$D=4$ and $D=3$ space dimensions with momentum and three ratios of the vector repulsion to the
scalar attraction, with the Fermi momentum of the bound state at the threshold (units $m=1$). The
ratio 0.2025 is that of the Hartree-Fock couplings for $g_1=-0.1\,g_0$. Last column: the critical
coupling of the vacuum for a cutoff equal to $m$ (Sec.~\ref{sec:cjtcontact}).}
\begin{ruledtabular}
\begin{tabular}{ccccc}
$D$ & $\GV/|\GS|$ & $|\GS|\,m^{D-1}$ & $\kF/m$ & $|G_c|\,m^{D-1}$ \\
\hline
{{ROWS:tbl_ks_thresholds}}
\end{tabular}
\end{ruledtabular}
\end{table}

If the coupling at a reference time is $\GS^{\mathrm{ref}}$, the binding begins at the clock reading
$\td=\ln(|\GS|_{\mathrm{threshold}}/|\GS^{\mathrm{ref}}|)$. Figure~\ref{fig:ksbind} follows the saturated
matter from that moment on: its energy per particle falls and the mass of its particles falls, more
slowly if a vector repulsion is present. The thresholds were computed twice, by the program and by
the independent script, and agree to $10^{-9}$.

\begin{figure*}
\includegraphics[width=0.86\textwidth]{fig_ks_binding}
\caption{\label{fig:ksbind}Self-bound matter in the sector of fields that do not depend on
$x5,x6,x7$, along the deflation clock, counted from the moment at which matter first becomes
self-bound. Upper row: energy per particle of the saturated state minus $m$. Lower row: mass of its
particles. Left: $D=4$; right: $D=3$. Three ratios of the vector to the scalar coupling.}
\end{figure*}

Matter of a \emph{fixed} density is carried through the same sequence
(Fig.~\ref{fig:ksgas}). Its pressure along the flat coordinates becomes negative at a moment that
depends on the density (Table~\ref{tab:ksgas}): a uniform gas is then unstable against breaking up
into drops of saturated matter. Along the deflating coordinates the pressure is the interaction
pressure $\pint$ of Eq.~(\ref{eq:pKS}) alone, which is negative as long as the scalar attraction
dominates: \emph{the matter pulls on the deflating coordinates.} Section~\ref{sec:gravity} uses this.

\begin{table}
\caption{\label{tab:ksgas}A gas of fixed Fermi momentum in the sector of fields that do not depend on
$x5,x6,x7$: the coupling at which its pressure along the flat coordinates vanishes, and the clock
reading of that moment counted from the binding threshold of Table~\ref{tab:ksthr} (negative: before
the threshold is reached, at this higher density).}
\begin{ruledtabular}
\begin{tabular}{ccccc}
$D$ & $\GV/|\GS|$ & $\kF/m$ & $|\GS|\,m^{D-1}$ & $\td$ \\
\hline
{{ROWS:tbl_ks_gas}}
\end{tabular}
\end{ruledtabular}
\end{table}

\begin{figure*}
\includegraphics[width=\textwidth]{fig_ks_gas}
\caption{\label{fig:ksgas}A gas of fixed density in the sector of fields that do not depend on
$x5,x6,x7$ ($D=4$, $\GV=0.1\,|\GS|$), along the deflation clock counted from the binding threshold.
(a) Mass of the particles in matter. (b) Pressure along $x0,\dots,x3$ divided by the energy density.
(c) Pressure along $x5,x6,x7$, the interaction pressure alone, divided by the energy density.}
\end{figure*}

%%<1644>
\subsection{The two halves of the field in the covariant form}
\label{sec:kshalves}

With the covariant form of the self-interaction and the positive kinetic term
(Sec.~\ref{sec:forms}) the particles of the two halves $W_{1644}=\pm1$ of the field are two kinds of
particles with four internal states each, which carry opposite scalar and vector charges. The
interaction energy is the four-term expression of Sec.~\ref{sec:hf}, and the masses of the two kinds
obey two coupled gap equations,
\begin{equation}
m^{*}_{\pm}=m+\GS^{\mathrm c}\,\nS\pm\GS^{\mathrm{HF}}\,n_{SW} .
\end{equation}
For a scalar attraction $\GS^{\mathrm{HF}}<0$ a particle is attracted by particles of its own half
and repelled by those of the other. The state with equal halves is then never the state of lowest
energy at a given total density: \emph{the matter polarizes spontaneously}
(Fig.~\ref{fig:kshalves}, Table~\ref{tab:kshalves}). The gap equations have exactly one solution at
each polarization, with positive masses (the independent script searched from 81 starting values,
including negative masses). For couplings three times the binding threshold of the quantized-sector
form and a total density of the Fermi momentum $0.8\,m$ the energy per particle minus $m$ is
${{GENM:halves_e0}}\,m$ for equal halves and ${{GENM:halves_e1}}\,m$ with all particles in one half,
against ${{GENM:halves_eq}}\,m$ in the quantized-sector form: \emph{the covariant form binds matter
much less.} A fully polarized state is matter of one kind of particle with four internal states.

\begin{table}
\caption{\label{tab:kshalves}Covariant form on $4+4$: the polarization
$\mathrm{pol}_W=(n_+-n_-)/n$ of the state of lowest energy, for couplings $g_0<0$, $g_1=-0.1\,g_0$
of 1.5, 3 and 6 times the value at which the quantized-sector form first binds matter, and four
total densities given by the Fermi momentum that the density would have with equal halves ($D=4$;
steps of 0.05 in the polarization).}
\begin{ruledtabular}
\begin{tabular}{ccccc}
couplings & $\kF=0.4\,m$ & $0.6\,m$ & $0.8\,m$ & $1.0\,m$ \\
\hline
{{ROWS:tbl_ks_halves}}
\end{tabular}
\end{ruledtabular}
\end{table}

\begin{figure*}
\includegraphics[width=\textwidth]{fig_ks_halves}
\caption{\label{fig:kshalves}Covariant form of the self-interaction on $4+4$: energy per particle as
a function of the polarization of the two halves of the field, relative to the state with equal
halves, for three strengths of the couplings and four total densities. Circles mark the state of
lowest energy.}
\end{figure*}
%%</1644>
